% ===================================================================
%  جدول نمبر 2 — انسانی جسم۔ تفریح کا وقت۔ کھیل۔
%
%  Traduction, faite en 2026, en ourdou d'aujourd'hui, sur l'ido de
%  texto/io. Regles de la colonne : voir 10-jadval-01.tex. Ecrit avec
%  outils/ordo.py sous les yeux.
%
%  L'ANATOMIE EST LE MEILLEUR TEST QU'UNE LANGUE PUISSE SUBIR, et
%  celle-ci le passe d'une facon qu'aucune autre colonne du releve
%  n'a montree : elle le passe DEUX FOIS. Presque chaque partie du
%  corps a en ourdou un mot indien et un mot arabe, et les deux
%  s'impriment — ماتھا et پیشانی pour le front, آنکھ et چشم pour
%  l'œil, دل et قلب pour le cœur, ہڈی et عظم pour l'os, ناک et بینی
%  pour le nez.
%
%  LA COLONNE CHOISIT LE MOT INDIEN, ET IL FAUT DIRE POURQUOI. Ce
%  n'est pas par purisme a rebours : c'est que ce livret est un
%  manuel de vocabulaire pour enfants, et que le mot du corps qu'un
%  enfant entend est ماتھا, pas پیشانی. Le second appartient a la
%  poesie et a la medecine. Choisir le registre, ici, c'est choisir
%  le lecteur — et le lecteur du fac-simile est un ecolier.
%
%  MAIS DEUX MOTS ARABES RESTENT, PARCE QU'ILS N'ONT PAS DE
%  CONCURRENT : عضلہ le muscle et شریان l'artere. L'ourdou courant
%  n'a rien d'autre, et forger un synonyme indien pour l'occasion
%  serait exactement la faute inverse.
%
%  LES CINQ SENS SONT ARABES TOUS LES CINQ — بینائی, سماعت, شامہ,
%  ذائقہ, لمس —, et la aussi c'est l'usage qui tranche : ce sont les
%  mots de l'ecole, ceux qu'on apprend en liste. Le corps se dit en
%  indien, ses facultes en arabe, et la frontiere passe exactement la
%  ou elle passait au tableau 1 entre le mobilier et les sciences.
%
%  LES JEUX SE PARTAGENT EN TROIS, COMME PARTOUT :
%
%   * ceux que le pays joue et nomme : لٹو la toupie (18), کنچے les
%     billes (21), گلی ڈنڈا — non, celui-la ne sert pas —, پتنگ le
%     cerf-volant (33), جھولا la balancoire (35), آنکھ مچولی le
%     colin-maillard (31), چھپن چھپائی cache-cache (25) ;
%   * ceux qui sont venus avec leur nom : ٹینس (34), کروکے (15),
%     بلبوکے (19) ;
%   * ceux qu'il faut decrire : le jeu de barres (38), la main
%     chaude (28), le jeu de l'ours (32), les quilles (16).
%
%  ET LE PIEGE DU JEU DE BARRES REVIENT UNE TROISIEME FOIS. Le coreen
%  avait 진놀이 sous la main, le tamoul கபடி ; l'ourdou a کبڈی, le
%  meme jeu sous le meme nom, et en outre le jeu national du
%  Pakistan. Trois langues, trois refus. La colonne ecrit روک ٹوک کا
%  کھیل, le jeu ou l'on barre le passage.
%
%  LE SAUT-MOUTON ET SA NOTE. Le fac-simile distingue le
%  shultro-salto (30), ou les sauteurs appuient les mains sur les
%  EPAULES, et le dorso-salto, ou ils n'appuient que sur le DOS. La
%  colonne coreenne avait d'abord interverti les deux ; celle-ci
%  ecrit کندھا پھلانگ et پیٹھ پھلانگ, dans cet ordre.
%
%  UN SEUL RETOURNEMENT SUR LES SEPT PAIRES DE parigi.py :
%  « kegli (16) quin li abatas per bulo (17) », la relative que
%  l'ourdou ne peut pas postposer. Meme compte qu'au coreen et au
%  tamoul, au meme tableau : le corps puis les jeux, et l'enumeration
%  ne s'inverse pas.
% ===================================================================

%%K t02-apar-1 apar
\VUsaut{4.0mm}
\VUtitre{15.0pt}{60}{جدول نمبر 2}
\VUsaut{2.0mm}
\VUtitre{11.4pt}{0}{انسانی جسم۔ --- تفریح کا وقت۔}
\VUtitre{11.4pt}{0}{کھیل۔}

%%K t02-tit-0 sub
\VUtitre{13.2pt}{0}{انسانی جسم۔}
\VUsaut{2.4mm}
\VUcentre{10.2pt}{0}{\textit{(علومِ فطرت کا ایک سادہ سبق۔)}}
\VUsaut{3.0mm}

%%K t02-tit-1 sub pos=t02-01-1
\VUpk{10.2pt}{40}{سر۔}
\VUsaut{3.0mm}

%%K t02-01-1 p
1. --- انسانی جسم کے بڑے حصے یہ ہیں : \VUgras{سر}، \VUgras{دھڑ} اور
\VUgras{اعضا}۔

%%K t02-02-1 p
2. --- سر میں دو حصے ہیں : \VUgras{کھوپڑی}، جس کے اندر \VUgras{دماغ}
ہے اور جسے \VUgras{بال} ڈھانپتے ہیں ؛ اور \VUgras{چہرہ}۔

%%K t02-03-1 p
3. --- ہماری تصویر میں نہ کھوپڑی نظر آتی ہے، نہ دماغ ؛ مگر چہرے کے
اہم حصے ہمیں بہت صاف نظر آتے ہیں۔

%%K t02-04-1 p
4. --- سب سے پہلے یہ \VUgras{ماتھا} ہے ؛ یہ زور دار عقل اور ہمیشہ
کام کرنے والی سوچ کا پتا دیتا ہے۔ \VUgras{کنپٹیاں} بہت کشادہ ہیں۔
\VUgras{آنکھ} چمکدار اور خوب کھلی ہوئی ہے۔ گھنی \VUgras{بھویں} اور
لمبی \VUgras{پلکیں} اُس کی حفاظت کرتی ہیں۔

%%K t02-05-1 p
5. --- یہ رہے \VUgras{ناک} اور \VUgras{نتھنے}۔ ناک سے ہم سونگھتے اور
سانس لیتے ہیں۔ ناک کے دونوں طرف \VUgras{گال} ہیں، اور کچھ دور
\VUgras{کان}۔ چہرے کا نچلا حصہ \VUgras{منہ} اور \VUgras{ٹھوڑی} سے
بنتا ہے۔

%%K t02-06-1 p
6. --- ہمیں وہ اچھی طرح نظر نہیں آتے، کیونکہ \VUgras{مونچھیں} اور
\VUgras{گلمچھے} اُنھیں ہم سے چھپا دیتے ہیں۔ مگر ہم جانتے ہیں کہ منہ
میں دو \VUgras{ہونٹ}، \VUgras{اوپر کا جبڑا} اور \VUgras{نیچے کا جبڑا}
اپنے \VUgras{دانتوں} کے ساتھ، اور \VUgras{زبان} ہوتی ہے۔ منہ سے ہم
بولتے اور کھاتے ہیں۔ زبان سے کھانے کا مزہ چکھتے ہیں۔

%%K t02-07-1 p
7. --- آنکھ، کان، ناک اور منہ، انگلیوں کی طرح، پانچ حواس کے اعضا
ہیں : \VUgras{بینائی}، \VUgras{سماعت}، \VUgras{شامہ}، \VUgras{ذائقہ}،
\VUgras{لمس}۔

%%K t02-08-1 p
8. --- پس سر انسانی جسم کے سب سے دلچسپ حصوں میں سے ایک ہے۔

%%K t02-tit-2 sub
\VUsaut{4.4mm}
\VUpk{10.2pt}{40}{دھڑ۔}
\VUsaut{3.0mm}

%%K t02-09-1 p
9. --- \VUgras{گردن} سر کو دھڑ سے ملاتی ہے۔ اس میں \VUgras{گلا} اور
\VUgras{گدی} شامل ہیں۔

%%K t02-10-1 p
10. --- دھڑ میں کئی ضروری اعضا بھی ہیں۔ \VUgras{سینے} میں
\VUgras{پھیپھڑے} ہیں، جن سے ہم سانس لیتے ہیں، اور \VUgras{دل} ہے، جو
زندگی کا مرکز ہے۔ جب دل دھڑکنا بند کر دے تو جاندار مر جاتا ہے۔

%%K t02-11-1 p
11. --- \VUgras{پسلیاں} سینے کو سنبھالتی ہیں۔

%%K t02-12-1 p
12. --- دھڑ کے باقی حصے \VUgras{پہلو}، \VUgras{پیٹھ}، \VUgras{کندھے}
اور \VUgras{پیٹ} ہیں۔ ہم سونے کے لیے پہلو یا پیٹھ کے بل لیٹتے ہیں ؛
بوجھ اپنے کندھے پر اٹھاتے ہیں۔

%%K t02-13-1 p
13. --- پیٹ کے اندر \VUgras{معدہ} اور \VUgras{آنتیں} ہیں۔ وہ ہمیں
کھانا ہضم کرنے میں مدد دیتی ہیں۔

%%K t02-tit-3 sub
\VUsaut{4.4mm}
\VUpk{10.2pt}{40}{اعضا۔}
\VUsaut{3.0mm}

%%K t02-14-1 p
14. --- ہمارے چار \VUgras{اعضا} ہیں : دو \VUgras{بازو} اور دو
\VUgras{ٹانگیں}۔ بازوؤں سے ہم چیزیں پکڑتے، تھامتے، اٹھاتے اور جمع
کرتے ہیں ؛ ٹانگوں سے کھڑے ہوتے، چلتے اور دوڑتے ہیں۔

%%K t02-15-1 p
15. --- \VUgras{کندھا} بازو کو جسم سے ملاتا ہے۔ ایک بہت مضبوط
\VUgras{عضلہ}، بائسپس، بازو کو ہلاتا ہے ؛ اُس بازو کو \VUgras{کہنی}
\VUgras{کلائی کے اوپر والے حصے} سے ملاتی ہے۔ \VUgras{گٹّا}
\VUgras{ہاتھ} کا جوڑ بناتا ہے ؛ ہاتھ \VUgras{انگلیوں} اور
\VUgras{ناخنوں} پر ختم ہوتا ہے۔ ہاتھ پر ہم \VUgras{رگ} (اور شریان)
پہچانتے ہیں ؛ اُن میں \VUgras{خون} بہتا ہے۔

%%K t02-16-1 p
16. --- ٹانگ کے حصے یہ ہیں : \VUgras{ران}، \VUgras{گھٹنا}،
\VUgras{گھٹنے کی پشت}، \VUgras{پنڈلی} — جس کے \VUgras{گوشت} کو
\VUgras{کھال} ڈھانپتی ہے —، \VUgras{ٹخنہ} اور \VUgras{پاؤں}۔ ٹانگ کی
\VUgras{ہڈیاں} بہت موٹی اور بہت مضبوط ہیں۔ پاؤں کے سرے پر
\VUgras{پاؤں کی انگلیاں} ہیں۔ باقی حصے \VUgras{تلوا} اور
\VUgras{ایڑی} ہیں۔

%%K t02-17-1 p
17. --- انسانی جسم کی تقریباً مکمل وضاحت یہی ہے۔

%%K t02-17-2 p
\textit{نوٹ۔} --- \textit{« مجھے... (سر وغیرہ) میں درد ہے » اس}
\textit{جملے کی مدد سے استاد یہ سارے الفاظ اچھی یا بری}
\VUgras{جسمانی حالت} \textit{کے پہلو سے دوبارہ استعمال کرا سکتا ہے۔}

%%K t02-tit-4 sub
\VUsaut{5.0mm}
\VUpk{10.2pt}{40}{ورزش۔}
\VUsaut{2.4mm}
\VUcentre{10.2pt}{0}{\textit{(ایک طالب علم کی زبانی۔)}}
\VUsaut{3.0mm}

%%K t02-18-1 p
18. --- لچکدار، پھرتیلا اور تندرست رہنے کے لیے ہمیں اپنی ٹانگوں اور
بازوؤں کو ورزش کرانی چاہیے۔ اسی لیے ہم کھیلتے ہیں اور \VUgras{ورزش}
کرتے ہیں۔

%%K t02-19-1 p
19. --- ہفتے کے دوران یہ دونوں مشقیں ثانوی اسکول کے صحن میں ہوتی ہیں
(جدول نمبر 1 دیکھیے)۔ مگر جمعرات اور اتوار کو ہم ایک بڑے سبزہ زار
میں جاتے ہیں ؛ وہاں دوڑنے اور کھیلنے کی جگہ ہوتی ہے۔

%%K t02-20-1 p
20. --- ہمارے استاد اور ہمارے والدین کبھی کبھی ہمارے ساتھ وہاں آتے
ہیں ؛ مگر مشقیں چلاتے تو \VUgras{ورزش کے
استاد}\textsuperscript{(12)} ہی ہیں۔

%%K t02-21-1 p
21. --- سبزہ زار میں، داخلی دروازے کی بائیں طرف، \VUgras{ورزشی
آلات}\textsuperscript{(1)} ہیں۔ \VUgras{سیڑھی}\textsuperscript{(2)} سے
ہم چڑھتے ہیں ؛ \VUgras{حلقوں}\textsuperscript{(3)} اور \VUgras{لٹکتی
ڈنڈی}\textsuperscript{(4)} پر الٹے لٹکتے ہیں ؛ \VUgras{رسے کی
سیڑھی}\textsuperscript{(5)} یا \VUgras{گرہ دار
رسے}\textsuperscript{(6)} کی چوٹی تک ہاتھوں سے اپنے آپ کو کھینچتے
ہیں۔

%%K t02-22-1 p
22. --- اُسی وقت ہمارے ساتھی \VUgras{لکڑی کے
گھوڑے}\textsuperscript{(7)} یا \VUgras{متوازی ڈنڈوں}\textsuperscript{(11)}
کے اوپر سے کودتے ہیں۔ کچھ اور \VUgras{مکے بازی
کرتے}\textsuperscript{(8)} ہیں یا منہ پر جالی پہن کر \VUgras{پتلی تلوار}
سے \VUgras{تلوار بازی کرتے}\textsuperscript{(9)} ہیں۔ آخر میں کچھ اور
\VUgras{وزن اٹھاتے}\textsuperscript{(10)} ہیں۔

%%K t02-tit-5 sub
\VUsaut{4.6mm}
\VUpk{10.2pt}{40}{کھیل۔}
\VUsaut{3.0mm}

%%K t02-23-1 p
23. --- ورزش کرنے والوں کے گرد لڑکے اور لڑکیاں کئی \VUgras{کھیل}
کھیلتے ہیں۔ لڑکیاں اپنے \VUgras{بڑے چھلے}\textsuperscript{(14)} سے
کھیلتی ہیں ؛ \VUgras{رسّی}\textsuperscript{(13)} کودتی ہیں، یا اپنے
بھائیوں اور کزنوں کو \VUgras{کروکے}\textsuperscript{(15)} کے کھیل میں
بلاتی ہیں۔ جس لمحے مخالف سمجھتا ہے کہ وہ محراب کے نیچے سے آسانی سے
گزر جائے گا، اُسی لمحے اُس کی \VUgras{گیند} کو اپنے \VUgras{لکڑی کے
ہتھوڑے} سے پرے دھکیلنے میں اُنھیں مزہ آتا ہے!

%%K t02-24-1 p
24. --- لڑکے زیادہ زور والے کھیل پسند کرتے ہیں۔ وہ شوق سے
\VUgras{ڈنڈوں}\textsuperscript{(16)} سے کھیلتے ہیں ؛ اُنھیں مہارت سے
\VUgras{گیند}\textsuperscript{(17)} سے گراتے ہیں۔ \VUgras{لٹو}
\textsuperscript{(18)}، \VUgras{بلبوکے}\textsuperscript{(19)}، حتیٰ کہ
\VUgras{مینڈک کے کھیل}\textsuperscript{(20)} کو بھی وہ حقیر نہیں
سمجھتے۔ مگر اُن کے سب سے پسندیدہ کھیل \VUgras{کنچے}\textsuperscript{(21)}
اور \VUgras{دیوار گیند}\textsuperscript{(22)} ہیں ؛ کیونکہ
\VUgras{شیشہ گھر}\textsuperscript{(26)} کی دیوار ہلکی گیند کو واپس
اچھالنے کے لیے بہت موزوں ہے۔

%%K t02-25-1 p
25. --- ننھا یوآنیس اپنے \VUgras{لمبے پایوں}\textsuperscript{(24)} پر
چلتا ہے ؛ وہ بہت ماہر ہے اور توازن رکھنا خوب جانتا ہے۔

%%K t02-noto-1 noto
\VUnotes{143.2mm}{%
(*) بچوں اور لڑکوں کے کھیلوں کے اکثر خالص قومی نام ہوتے ہیں۔ ہم نے
اُنھیں اِدو میں جہاں تک ہو سکا اچھی طرح نام دینے کی کوشش کی — کبھی
اُس کام سے جس سے وہ پہچانے جاتے ہیں، کبھی اِس یا اُس ملک کا نام چن
کر، جب وہ زیادہ نامناسب نہ ہو۔ مثلاً : \VUgras{پیپے کا کھیل}
(جرمنی اور فرانس میں) وہی \VUgras{مینڈک کا کھیل} \textsuperscript{(20)}
ہے (انگریزی میں) ؛ اُس میں کھیلنے والے اپنی گول ٹکیاں سوراخ دار
صندوق (پیپے) پر منہ کھولے کھڑے دھات کے مینڈک کے منہ میں پھینکنے کی
کوشش کرتے ہیں۔
\VUgras{کندھا پھلانگ}\textsuperscript{(30)} \VUgras{پیٹھ پھلانگ} سے
صرف اِس میں مختلف ہے : سر کے اوپر سے کودنے کے لیے کھیلنے والے اپنے
ہاتھ ساتھی کے \textit{کندھوں} پر ٹیکتے ہیں ؛ « \VUgras{پیٹھ
پھلانگ} » میں صرف اُس کی \textit{پیٹھ} پر ٹیکتے ہیں۔ --- یا پھر
شریک ہونے والوں میں سے کچھ جھک کر کھڑے ہوتے ہیں اور باقی ہاتھ کندھے
پر ٹیک کر اُن کی پیٹھ کے اوپر سے کودتے ہیں ؛ پہلوں کو جھٹکا بغیر
جھکے سہنا ہوتا ہے، دوسروں کو توازن کھوئے بغیر کودنا ہوتا ہے (خود
جدول دیکھیے)۔%
}

%%K t02-25-1 p suite
اُس کے قریب کچھ ساتھی اُسی شیشہ گھر میں اور \VUgras{باڑ} کے پیچھے
\VUgras{چھپن چھپائی}\textsuperscript{(25)} کھیلتے ہیں۔ کچھ اور
\VUgras{مارنے والے کو پہچاننا}\textsuperscript{(28)} یا ایک
\VUgras{بلّے} سے دوسرے بلّے کی طرف اڑنے والی \VUgras{اڑتی
گیند}\textsuperscript{(29)} پسند کرتے ہیں۔

%%K t02-26-1 p
26. --- سبزہ زار کے بیچ میں دو درخت ہیں۔ اُن میں سے ایک کے نیچے سات
آٹھ لڑکوں نے \VUgras{کندھا پھلانگ}\textsuperscript{(30)} (*) کا کھیل
جمایا ہے۔ ہر ملک میں یہ کھیل نہیں جانا جاتا ؛ مگر \VUgras{پیٹھ
پھلانگ} کیا ہے، یہ سب جانتے ہیں — اِس بار لڑکے وہ نہیں کھیل رہے۔

%%K t02-tit-6 sub
\VUsaut{5.0mm}
\VUpk{10.2pt}{40}{کھلی ہوا کے کھیل۔}
\VUsaut{3.4mm}

%%K t02-27-1 p
27. \VUgras{آنکھ مچولی}\textsuperscript{(31)} بھی بہت مزے کا کھیل ہے ؛
لڑکیاں اور لڑکے باری باری اپنی آنکھوں پر \VUgras{پٹی} باندھتے ہیں اور
اُن ناتجربہ کاروں کو پکڑتے ہیں جو خود کو پکڑوا لیتے ہیں۔

%%K t02-28-1 p
28. --- سبزہ زار کے بیچ میں، یہ رہا ایک بہت فرانسیسی مگر کچھ سخت
کھیل : وہ \VUgras{ریچھ کا کھیل}\textsuperscript{(32)} ہے ؛ یہ خاموش
\VUgras{پتنگ}\textsuperscript{(33)} کے کھیل سے، حتیٰ کہ انگریزی کھیل
\VUgras{ٹینس}\textsuperscript{(34)} سے بھی کہیں زیادہ شور والا ہے۔

%%K t02-29-1 p
29. --- یہ آخری کھیل بڑے طالب علموں کی خوشی ہے۔ ہمارے استاد خود شوق
سے کھیلتے ہیں۔ اِس کے لیے بڑی مہارت اور پھرتی چاہیے، اور مجھے یہ بہت
پسند ہے۔ \VUgras{جھولے}\textsuperscript{(35)} کے کھیل میں لڑکیوں ہی کے
لیے چھوڑ دیتا ہوں۔

%%K t02-30-1 p
30. --- ایک اور انگریزی کھیل مجھے پسند نہیں، جو شاید خطرناک ہو جائے۔

%%K t02-30-1b p
میں \VUgras{فٹ بال}\textsuperscript{(36)} کی بات کر رہا ہوں ؛ وہ میرے
بڑے بھائی کو خوش کرتا ہے۔ میں اپنا پرانا \VUgras{روک ٹوک کا
کھیل}\textsuperscript{(38)} زیادہ پسند کرتا ہوں ؛ وہ اُتنا ہی صحت
بخش اور سچ مچ کم سخت ہے۔

%%K t02-tit-7 sub
\VUsaut{4.6mm}
\VUpk{10.2pt}{40}{سائیکل اور گیند۔}
\VUsaut{3.0mm}

%%K t02-31-1 p
31. --- \VUgras{سائیکل}\textsuperscript{(37)} پر سبزہ زار کے گرد چکر
لگانا بھی مجھے پسند ہے۔

%%K t02-32-1 p
32. --- اگلے سال \VUgras{گھڑ سواری}\textsuperscript{(39)} سیکھوں گا۔
خوبصورتی سے قدم اٹھاتا، تیز چلتا اور رکاوٹوں کے اوپر سے چھلانگ لگاتا
گھوڑا کتنا خوبصورت ہوتا ہے!

%%K t02-33-1 p
33. --- گرمیوں میں ہم \VUgras{تیراکی}\textsuperscript{(40)} سیکھتے
ہیں۔ دریا کے صاف ٹھنڈے پانی میں مزے سے غوطہ لگاتے اور نہاتے ہیں۔
کبھی کبھی آخر میں ایک کاغذی \VUgras{غبارہ}\textsuperscript{(41)}
چھوڑتے ہیں ؛ گرم ہوا سے بھرتے ہی وہ شان سے فضا میں چڑھ جاتا ہے۔

%%K t02-34-1 p
34. --- اور بھی بہت سے کھیل ہیں، مگر اُنھیں گنتے گنتے میں ختم نہ کر
سکوں گا ؛ اِس خلاصے ہی سے سمجھ آ جاتا ہے کہ جمعرات کو ہمارے خوبصورت
سبزہ زار میں زیادہ اکتاہٹ نہیں ہوتی۔

%%K t02-34-2 p
تنبیہ۔ --- \textit{استاد اِس باب کو آسانی سے مکمل کر سکتا ہے ؛ سب سے}
\textit{اہم کھیلوں میں سے ہر ایک کی زیادہ تفصیل سے وضاحت کر کے۔}
